\documentclass[11pt]{article}
% ======================================================================
%  weekpreamble.tex — shared preamble for the weekly course summaries (EN)
%  Concise, readable, intuition-first. Same box style as the main doc.
%  Compiles with pdflatex.
% ======================================================================
\usepackage[utf8]{inputenc}
\usepackage[T1]{fontenc}
\usepackage[english]{babel}
\usepackage[a4paper,margin=2.2cm]{geometry}
\usepackage{amsmath,amssymb,amsthm,mathtools}
\usepackage{bm}
\usepackage{enumitem}
\usepackage{xcolor}
\usepackage[most]{tcolorbox}
\usepackage{microtype}

\emergencystretch=3em
\allowdisplaybreaks[1]
\setlength{\parindent}{0pt}
\setlength{\parskip}{0.45em}

% ---------- math macros (same names as the main document) ----------
\newcommand{\E}{\mathbb{E}}
\DeclareMathOperator{\Var}{Var}
\DeclareMathOperator{\Cov}{Cov}
\DeclareMathOperator{\Corr}{Corr}
\DeclareMathOperator*{\argmin}{arg\,min}
\DeclareMathOperator{\MSE}{MSE}
\DeclareMathOperator{\trace}{tr}
\newcommand{\Reals}{\mathbb{R}}
\newcommand{\Ints}{\mathbb{Z}}
\newcommand{\Comp}{\mathbb{C}}
\newcommand{\Nat}{\mathbb{N}}
\newcommand{\Prob}{\mathbb{P}}
\newcommand{\Tr}{^{\mathsf{T}}}
\newcommand{\Herm}{^{\mathsf{H}}}
\newcommand{\conj}[1]{\overline{#1}}
\newcommand{\abs}[1]{\left\lvert #1 \right\rvert}
\newcommand{\norm}[1]{\left\lVert #1 \right\rVert}
\newcommand{\ind}{\mathbf{1}}
\newcommand{\eps}{\varepsilon}
\newcommand{\diff}{\nabla}
\newcommand{\acvs}{\gamma}
\newcommand{\acf}{\rho}
\newcommand{\pacf}{\alpha}
\newcommand{\sdf}{\mathcal{S}}
\newcommand{\filt}{\mathcal{L}}
\newcommand{\Bsh}{B}
\newcommand{\ms}{\;\overset{\mathrm{ms}}{=}\;}
\newcommand{\toprob}{\xrightarrow{P}}
\newcommand{\todist}{\xrightarrow{d}}
\newcommand{\iid}{\overset{\text{iid}}{\sim}}

\setlist[itemize]{leftmargin=1.5em,topsep=2pt,itemsep=2pt}
\setlist[enumerate]{leftmargin=1.7em,topsep=2pt,itemsep=2pt}

% ---------- compact, titled (unnumbered) boxes ----------
\tcbset{wkbase/.style={enhanced,breakable,fonttitle=\bfseries,coltitle=white,
  boxrule=0.6pt,arc=1.2mm,left=2.2mm,right=2.2mm,top=1.1mm,bottom=1.1mm,
  before skip=6pt,after skip=6pt,
  attach boxed title to top left={yshift=-3.2mm,xshift=2.4mm},
  boxed title style={boxrule=0pt,arc=0.5mm}}}

\newtcolorbox{defn}[1]{wkbase, colback=blue!4!white, colframe=blue!58!black,
  colbacktitle=blue!58!black, title={Definition --- #1}, top=3mm}
\newtcolorbox{result}[1]{wkbase, colback=red!4!white, colframe=red!60!black,
  colbacktitle=red!60!black, title={#1}, top=3mm}
\newtcolorbox{intu}[1][Intuition]{wkbase, colback=yellow!12!white,
  colframe=orange!72!black, colbacktitle=orange!72!black, coltitle=white,
  title={#1}, top=3mm}
\newtcolorbox{retenir}[1][Key takeaways --- the week's reflexes]{wkbase,
  colback=green!5!white, colframe=green!48!black, colbacktitle=green!48!black,
  title={#1}, top=3mm}
\newtcolorbox{exmpl}[1][Worked example]{wkbase, colback=black!3!white,
  colframe=black!55!white, colbacktitle=black!55!white, title={#1}, top=3mm}

% ---------- week header ----------
\newcommand{\weekheader}[4]{%
  {\small\sffamily\bfseries\color{black!60} TIME SERIES~~$\cdot$~~MATH-342, EPFL}\par
  \vspace{2pt}
  {\LARGE\bfseries Week #1 --- #3\par}
  \vspace{1pt}
  {\itshape\color{black!55} #2}\par
  \vspace{6pt}
  \begin{tcolorbox}[enhanced,colback=blue!3!white,colframe=blue!45!black,
    arc=1.4mm,boxrule=0.7pt,left=3mm,right=3mm,top=2mm,bottom=2mm,before skip=2pt,after skip=10pt]
    \textbf{The big picture.}~ #4
  \end{tcolorbox}
}

\providecommand{\Dir}{\mathcal{D}}
\begin{document}
\weekheader{06}{26 March 2026}{Fourier Theory}{The Fourier transform (continuous, discrete, finite-data) and the convolution theorem move us from time to frequency; beware aliasing.}

The Fourier transform is the bridge between the \textbf{time domain} (a function or a sequence $g_t$) and the \textbf{frequency domain} (its transform $G(f)$). We build it in three stages: continuous time (the whole function), discrete time (a sampled sequence), then \textbf{finite data} (what we actually observe). Throughout we use ordinary frequency $f$ (not angular frequency $\omega=2\pi f$). This machinery is the toolbox we will reuse next week to define the spectral density.

\section*{Key definitions}

\begin{defn}{Continuous-time Fourier transform}
For $g:\Reals\to\Comp$ with $g\in L^1$,
\[
G(f)=\int_{-\infty}^{\infty} g(t)\,e^{-2\pi i f t}\,dt,\qquad f\in\Reals .
\]
If moreover $G\in L^1$, the \textbf{inverse} recovers $g$ for almost all $t$: $g(t)=\int_{-\infty}^{\infty} G(f)\,e^{2\pi i f t}\,df$. We write $g\longleftrightarrow G$ (a \textbf{Fourier pair}): no information is lost. Upper case denotes the transform.
\end{defn}

\begin{defn}{Convolution}
For $g,h\in L^1$, $u=g*h$ is $u(t)=\int_{-\infty}^{\infty} g(t-u)\,h(u)\,du$. It is a \textbf{weighted average} of one function by the other, centred at $t$.
\end{defn}

\begin{defn}{Discrete-time Fourier transform}
For a sequence $\{g_t\}_{t\in\Delta\Ints}\in\ell^1$,
\[
G(f)=\Delta\sum_{t\in\Delta\Ints} g_t\,e^{-2\pi i t f},\qquad
g_t=\int_{-1/(2\Delta)}^{1/(2\Delta)} G(f)\,e^{2\pi i t f}\,df .
\]
The inverse is just a Fourier series with the domains swapped. Key feature: $G$ is \textbf{periodic}, $G(f)=G(f+k/\Delta)$, so it is fully determined on one period of length $1/\Delta$.
\end{defn}

\begin{defn}{Finite-data transform and the Dirichlet kernel}
With $n$ observations at times $T=\{0,\Delta,\dots,(n-1)\Delta\}$,
\[
G_n(f)=\sum_{t\in T} g_t\,e^{-2\pi i t f}.
\]
$G_n$ is the \emph{full} transform times a window indicator $d_t=\tfrac1\Delta\ind_T(t)$, whose transform is the \textbf{Dirichlet kernel}
\[
\Dir_T(f)=\sum_{t\in T}e^{-2\pi i f t}=\frac{\sin(n\pi f\Delta)}{\sin(\pi f\Delta)}\,e^{-\pi i(n-1)f\Delta}.
\]
Hence $G_n=\Dir_T * G$: a finite window \textbf{blurs} the true transform.
\end{defn}

\begin{defn}{Fourier frequencies}
The $n$ equally spaced frequencies (over one period) at which the FFT evaluates $G_n$:
\[
f_k=\frac{k}{\Delta n},\qquad -\frac{n}{2}\le k\le \frac{n-1}{2}
\quad\Big(=\tfrac{k}{n}\text{ if }\Delta=1\Big).
\]
\end{defn}

\section*{Key results \& formulas}

\begin{result}{Theorem --- Fourier inversion}
If $g,G\in L^1$ (continuous time) or $\{g_t\}\in\ell^1$ (discrete time), the inverse transform recovers $g$. \emph{Idea:} substitute the forward transform and use the orthogonality relation $\int e^{2\pi i f(t-s)}\,df$, which acts like a delta. This is what makes $(g,G)$ a genuine pair and underlies every time$\leftrightarrow$frequency duality.
\end{result}

\begin{result}{Theorem --- Convolution theorem}
\textbf{Multiplication in one domain = convolution in the other.} For $g,h\in L^1$:
\[
\begin{aligned}
u=g*h &\;\longleftrightarrow\; U=G\cdot H,\\
v=g\cdot h &\;\longleftrightarrow\; V=G*H \quad(\text{if }G,H\in L^1).
\end{aligned}
\]
In discrete time both pairs hold too: $h*g\leftrightarrow H\cdot G$ and $h\cdot g\leftrightarrow H*G$. \emph{Idea:} plug the convolution into the definition, swap the order of integration (Fubini, since $\norm{g}_1\norm{h}_1<\infty$), and recognise the phase factor. This is the central computational identity (filtering, blurring, etc.).
\end{result}

\begin{result}{Theorem --- Aliasing and the Nyquist frequency}
Sampling \emph{folds} all the continuous-frequency content onto one period:
\[
G(f)=\sum_{k\in\Ints} G\!\left(f+\tfrac{k}{\Delta}\right).
\]
The \textbf{Nyquist frequency} is $1/(2\Delta)$ ($=\tfrac12$ if $\Delta=1$): content above it is \textbf{not} recoverable. \emph{Idea:} write $g(t)$ via both inversion formulas at $t\in\Delta\Ints$, split $\Reals$ into blocks of length $1/\Delta$, shift each ($e^{2\pi i tk/\Delta}=1$ for these $t$), and match.
\end{result}

\begin{result}{Proposition --- Fast Fourier Transform (FFT)}
Evaluating $G_n$ at the $n$ Fourier frequencies costs $O(n\log n)$ (vs.\ $O(n^2)$ naively) by \textbf{divide-and-conquer} (even/odd indices, $\log_2 n$ levels). This makes large-scale spectral computation feasible.
\end{result}

\textbf{Basic properties} (same form in continuous and discrete time):
\[
\begin{aligned}
\conj{g(t)} &\longleftrightarrow \conj{G(-f)}, &\quad
g(\alpha t) &\longleftrightarrow \tfrac{1}{\abs\alpha}G(f/\alpha),\\
g(t+\tau) &\longleftrightarrow G(f)\,e^{2\pi i \tau f}, &\quad
\alpha g(t)+h(t) &\longleftrightarrow \alpha G(f)+H(f).
\end{aligned}
\]
(Conjugation reflects; stretching time compresses frequency; a time shift multiplies by a phase; the transform is linear.)

\section*{Intuition}

\begin{intu}[Convolution = product: the Swiss-army rule]
The rule ``\emph{convolution in one domain $\leftrightarrow$ product in the other}'' is the whole point of Fourier: an expensive operation (convolution) becomes a plain multiplication. That is why we go to the frequency domain to filter, and why the blurring caused by a finite window is written as a convolution with a kernel.
\end{intu}

\begin{intu}[Why aliasing hurts]
A \emph{high}-frequency sinusoid sampled too slowly looks \emph{exactly} like a low-frequency one --- the ``wagon-wheel'' effect in films. Frequencies $f$ and $f+k/\Delta$ become indistinguishable: they are ``aliases''. Cure: sample finely (small $\Delta$, large Nyquist) or low-pass filter \emph{before} sampling.
\end{intu}

\begin{intu}[Finite window = blurred vision]
We never observe the whole sequence, only $n$ points. Mathematically, observing = multiplying by an indicator. By the convolution theorem this \textbf{convolves} the true transform with the Dirichlet kernel $\Dir_T$: peaks get smeared and energy ``leaks'' into neighbouring frequencies. The larger $n$, the more concentrated $\Dir_T$, the less blurring.
\end{intu}

\begin{intu}[Discrete = continuous, but periodic and folded]
In continuous time $G$ lives on all of $\Reals$. As soon as we sample, $G$ becomes \emph{periodic} (period $1/\Delta$): knowing it on $[-\tfrac1{2\Delta},\tfrac1{2\Delta}]$ suffices. Aliasing says exactly how the continuous spectrum folds back into that period.
\end{intu}

\begin{exmpl}[Time shift]
Shifting a signal by $\tau$ does not change its amplitude content, only its phase: $g(t+\tau)\leftrightarrow G(f)e^{2\pi i\tau f}$, so $\abs{G}$ is unchanged. Intuition: moving a sound earlier or later in time does not change \emph{which} frequencies it contains.
\end{exmpl}

\section*{Key takeaways}

\begin{retenir}
\begin{itemize}
\item Be able to write the \textbf{three} transforms (continuous, discrete, finite-data) and their inverses, and to recognise a \textbf{Fourier pair} $g\leftrightarrow G$.
\item \textbf{Core reflex:} convolution $\leftrightarrow$ product. Use it to relate $G_n$ to $G$ via $G_n=\Dir_T*G$, and for any filtering computation.
\item Know the \textbf{4 basic properties} (conjugation, scaling, shift, linearity) --- one-line proofs by a change of variable / inserting a phase.
\item \textbf{Aliasing:} $G(f)=\sum_k G(f+k/\Delta)$; above the Nyquist frequency $1/(2\Delta)$ the information is lost. Sample fast enough.
\item In discrete time everything is \textbf{periodic} with period $1/\Delta$: work on one period only.
\item The \textbf{FFT} gives $G_n$ at the Fourier frequencies in $O(n\log n)$.
\item \emph{Next stop:} this machinery defines the spectral density $\sdf(f)=\sum_\tau\acvs_\tau e^{-2\pi i f\tau}$ --- the transform of the ACVS.
\end{itemize}
\end{retenir}

\end{document}
